% =============================================================================
% Chapter 5  The Descriptor Reformulation  --  ~9.5 pp  --  THE CORE
% rubric: Project Body (Modelling) -- Model Formulation /30.
% Plan: ThesisPlanning.md §5.5 (headline contribution C).
%
% PURPOSE     Answer Q(ii): the headline, written as the methods section of a
%             journal article, including the reformulation's own challenge and
%             its cure.
% CRITERION   Formulation, band 4, "warranting publication". The body carries
%             the idea, the object (stated once), the cost (a formal result),
%             the challenge and the interpretation; entries, constructions and
%             proofs go to Appendices B and G.
% SOURCES     ThesisPlanning.md §3.4, §3.7, §6.4; pihm/docs/02_methods.md §3.2;
%             pihm/docs/13_descriptor.md; pihm/docs/14_descriptor_gap.md (the
%             challenge); claims B-01-B-07, B-14, B-20, B-21, C-05, C-06;
%             pihm/docs/companion/{12,50-57}.
% WRITE WHEN  The mathematics now, with numbers as \res; in full at gate G4 (wave W3).
% EQUATION TEST (STYLE_GUIDE.md §2.6): does the argument break if a reader skips
%             it? If not, the equation belongs in an appendix.
%
% 2026-10-07 PROSE REWRITE (first pass for the author to rewrite; the previous
%   version is at commit e181c21). Changes: paragraph headings that announced a
%   register ("Cost, in the same breath", "Optimality, hedged", "What it measures",
%   "Interpretation") are replaced by content or removed; "the idea in one
%   sentence" is now the spine sentence itself; every benchmark problem is
%   introduced by its equation at first use, its tag after (author's choice);
%   NOTATION: the rescaling rule's characteristic scale no longer uses rho (rho is
%   the Carleman ratio, ThesisPlanning.md §8.2), and the band reflection is Q_m,
%   not R_m (R_0 is Chapter 4's row weight, R the stacked residual); the two side
%   points of the spectrum section (relative gap; one query's Haar success) are one
%   paragraph (author's choice); the per-query gate comparison is scoped to R2a (see
%   the REVISE note there). Kept from the previous draft: the duplicate label
%   eq:G-structured is eq:G-banded here; the chain's recurrence is stated in the raw
%   coefficients it is true for; the kernel theorem covers both layouts with the
%   proof in the body; the cost proposition's alpha/||A|| -> 1 is scoped and its
%   ancilla count is + 8; the atom table (tab:descriptor-atoms) stays dropped; the
%   CUT and MOVED notes below stand.
% =============================================================================

This chapter answers the second question of \Cref{sec:pihm-assessment}: can the residual be reformulated so that its Hamiltonian is banded and well conditioned without changing its solution, and at what cost? We build the reformulated residual, the descriptor form, and prove that its kernel is the standard form's (\Cref{sec:descriptor-idea,sec:descriptor-construction,sec:descriptor-layouts,sec:descriptor-constraints}). We then identify the difficulty it introduces, a spectral gap eroded by the carried derivatives, and the a priori rescaling that restores it on one axis (\Cref{sec:descriptor-gap,sec:descriptor-rescaling}). Finally we construct its block encoding and derive what it costs and what it does to the filter's degree (\Cref{sec:descriptor-block-encoding,sec:descriptor-cost,sec:descriptor-spectrum}).

\subsection{The Idea}
\label{sec:descriptor-idea}

\Cref{sec:standard} located the cost of the standard form in its Hamiltonian: $\|\G^k\|_2=\Theta(N^{2k})$, so $\|\Hstd\|=\Theta(N^{4k})$ whatever the encoder. That norm is the density of $\G$ made visible. The density is not intrinsic to differentiation; it comes from requiring the derivative's coefficients to be expressed in the same basis as the function's. If the derivative may instead live in a neighbouring polynomial basis, differentiation is banded, and $\G$ factors into two banded matrices whose product is dense only because one of them is inverted.

The descriptor form keeps the factors instead of the product. Where the standard form substitutes $\G^j$ for each derivative, the descriptor form carries $f',\dots,f^{(k)}$ as unknowns beside $f$, ties consecutive derivatives by banded recurrences, and writes the equation over all of them. A first-order equation or system needs no extra unknowns, because the equation itself gives the derivative: multiplying the standard rows through by one banded factor is enough. We call the first layout the \emph{chain} and the second the \emph{mass form} (\Cref{sec:descriptor-layouts}).

Neither ingredient is new to numerical analysis. Banded differentiation is the ultraspherical spectral method of Olver and Townsend \cite{olver2013}, of which the factors of \Cref{sec:descriptor-construction} are the discrete shadow (\Cref{sec:bg-ultraspherical}). Carrying derivatives as unknowns so as not to square the condition number of a least-squares functional is the idea of first-order-system least squares \cite{cai1994fosls} (\Cref{sec:bg-fosls}). The physics-informed Hamiltonian is a least-squares functional, so both apply to it. The contribution of this thesis is their transfer to a block-encoded physics-informed Hamiltonian: a proof that the kernel is unchanged, an encoding of the residual as one flat linear combination of unitaries, an a priori rescaling that restores the gap the new unknowns erode, and a measurement of what all of this does to ground-state preparation.

\Cref{fig:sparsity} shows the effect on the running example $f''+4f'+4f=0$ (R2a). At $n=7$ the standard form's stacked residual holds \res{build/R2a/standard/n7}{nonzeros} non-zero entries on \res{build/R2a/standard/n7}{dimension} unknowns, and the descriptor's holds \res{build/R2a/descriptor/n7}{nonzeros} on \res{build/R2a/descriptor/n7}{dimension}: four times as many unknowns, because of the block register that \Cref{sec:descriptor-layouts} counts, and six times fewer entries. That trade, a larger state for a sparser residual, is the subject of the chapter. Carrying the intermediate derivatives instead of eliminating them turns the residual from dense to banded with its kernel unchanged, so that the Hamiltonian's norm falls from $\Theta(N^{4k})$, which is $\Theta(N^8)$ at second order, to $\Theta(N^2)$ at fixed block weights; the price is a block register and a gap that the derivative blocks erode, which an a priori rescaling restores on one axis at no gate cost.

\placeholderfigure{\textbf{Dense against banded: the two residuals of the same equation.} Sparsity of the residual operators of $f''+4f'+4f=0$ (R2a) at $n=4$. Left, the standard form's $\Astd = \G^2 + 4\G + 4I$, whose upper triangle is full. Right, the descriptor's $\Adesc$, whose recurrence blocks, equation row and padding are banded. The descriptor's matrix is four times the size and, at $n=7$, holds six times fewer non-zero entries (\Cref{sec:descriptor-idea}).}{fig:sparsity}{Two sparsity (spy) plots side by side for $f''+4f'+4f=0$ (R2a) at $n = 4$: left, $\Astd = \G^2 + 4\G + 4I$ with its dense upper triangle; right, $\Adesc$ with its block-banded pattern (recurrence blocks, equation row, padding).}
% REVISE: "six times fewer" is 8384/1402 = 5.98 from the n7 keys above; if the
%   exporter's values move, recheck (or replace by a derived ratio key).

\subsection{The Banded Factorisation}
\label{sec:descriptor-construction}

Let $f=\sum_\ell a_\ell T_\ell$ have derivative $f'=\sum_\ell p_\ell T_\ell$. The two coefficient sequences are tied by the recurrence
\begin{equation}
    c_{\ell-1}\,p_{\ell-1} - p_{\ell+1} = 2\ell\,a_\ell ,
    \qquad \ell \ge 1,
    \label{eq:cheb-recurrence}
\end{equation}
with $c_0=2$ and $c_\ell=1$ for $\ell\ge1$, the factor on $p_0$ being an artefact of the normalisation of $T_0$ (the derivation is given in \Cref{app:cheb}). Each row ties one coefficient of the function to two of the derivative. Truncated to $N$ coefficients and written in the normalised basis of \Cref{sec:pihm-model}, with $\fhat{0}$ and $\fhat{1}$ the coefficients of $f$ and $f'$, it is the matrix identity
\begin{equation}
    \Bt\,\fhat{1} = \Dt\,\fhat{0} .
    \label{eq:banded-factorisation}
\end{equation}
$\Dt$ has a single non-zero diagonal, the first superdiagonal, with entries $(2\sqrt2,4,6,\dots,2(N-1))$; $\Bt$ has the diagonal $(2,1,1,\dots)$ and the second superdiagonal $(-\sqrt2,-1,-1,\dots)$. Both hold $\mathcal{O}(N)$ non-zero entries against the $\Theta(N^2)$ of $\G$. Two properties are used throughout: $\Bt$ is upper triangular with a non-zero diagonal, so it is invertible; and since the singular values of a single-diagonal matrix are its entries' magnitudes, $\|\Dt\|_2=2(N-1)$, while $\|\Bt\|_2\le 2+\sqrt2$.

The dense differentiation matrix is therefore the quotient
\begin{equation}
    \G = \Bt^{-1}\Dt ,
    \label{eq:G-banded}
\end{equation}
exactly at every $N$. The inverse $\Bt^{-1}$ is a triangle of ones on each parity class, and the density of $\G$ is that inverse; \Cref{eq:banded-factorisation} is the same linear map before the inverse is taken. This is the discrete form of the ultraspherical spectral method \cite{olver2013}: up to the basis normalisation, $\Bt$ and $\Dt$ are twice that method's operators for conversion from first-kind to second-kind Chebyshev polynomials and for differentiation into the second kind (\Cref{app:cheb}). The ultraspherical method applies these operators to eliminate the derivatives and obtain a banded system in $f$ alone. We instead keep \Cref{eq:banded-factorisation} as a row of the system and the derivative as an unknown, which lets one construction serve every order, every axis and every coupled field.

\subsection{Two Layouts and Kernel Equivalence}
\label{sec:descriptor-layouts}

Consider the linear \gls{ode} of order $k$
\begin{equation}
  \sum_{j=0}^{k} a_j(x)\, f^{(j)}(x) = r(x), \qquad x \in [-1,1],
\end{equation}
with the data constraints of \Cref{sec:pihm-residual}, every function in it represented by its first $N$ coefficients in the normalised Chebyshev basis. Which layout carries it depends on its order.

\paragraph{The chain ($k\ge2$).} Let $\fhat{j}\in\mathbb{C}^N$ be the coefficient vector of $f^{(j)}$ for $j=0,\dots,k$, stacked into the state
\begin{equation}
  w = (\fhat{0};\, \fhat{1};\, \dots;\, \fhat{k}).
\end{equation}
Two kinds of row couple the blocks. The \emph{recurrence rows}
\begin{equation}
  -\Dt\, \fhat{j} + \Bt\, \fhat{j+1} = 0, \qquad j = 0, \dots, k-1,
  \label{eq:descriptor-recurrence-rows}
\end{equation}
apply \Cref{eq:banded-factorisation} block by block. They contain no derivative operator, and on the kernel they are equivalent to $\fhat{j+1}=\G\fhat{j}$. The \emph{equation row}
\begin{equation}
  \sum_{j=0}^{k} \M_{a_j}\, \fhat{j} - \M_r\,\Dzero(x_s)\,\fhat{0} = 0
  \label{eq:descriptor-ode-row}
\end{equation}
is the equation itself, with $\M_{a_j}$ the multiplication operator of \Cref{sec:pihm-residual} (it reduces to $a_jI$ for a constant coefficient, and lifts the row to $n+1$ qubits wherever the standard form lifts it) and the source entering as it does there. The derivatives are unknowns of $w$ and are never replaced by $\G^j\fhat{0}$, so the dense quotient \Cref{eq:G-banded} never appears. The name is borrowed from descriptor systems in control theory, which likewise keep intermediate quantities as unknowns beside the state and tie them by algebraic rows.

\begin{definition}[Descriptor residual operator]
\label{def:descriptor-residual-operator}
Collecting the recurrence rows and the equation row, the descriptor residual operator $\Adesc$ is the block matrix
\begin{equation}
\Adesc =
\begin{pmatrix}
-\Dt & \Bt &        &        &          \\
          & -\Dt & \Bt &      &          \\
          &          & \ddots & \ddots &          \\
          &          &        & -\Dt & \Bt \\
\M_{a_0} - \M_r\Dzero & \M_{a_1}  & \cdots & \cdots & \M_{a_k}
\end{pmatrix}
\end{equation}
with $(k+1)N$ columns, so that $\Adesc\,w=0$ is the equation. Its first $k$ block rows are the recurrences and its last is the equation row.
\end{definition}

\paragraph{The padded block register.} On a quantum register the block index is held by $n_f=\lceil\log_2(k+1)\rceil$ qubits, which address $2^{n_f}$ blocks rather than $k+1$. The state is therefore padded to $2^{n_f}N$ amplitudes by blocks $b>k$, each pinned to zero by its own row $c_{\mathrm{pad}}\fhat{b}=0$ with $c_{\mathrm{pad}}=1$. No other row reads a padding block, so the padding decouples: the Hamiltonian is block diagonal, and the spectrum of the physical blocks is unaffected. The padding is inert but not free. The chain's state has $2^{n_f}N$ amplitudes against the standard form's $N$, which is $n_f$ more qubits and up to twice the unpadded $(k+1)N$ when $k+1$ is not a power of two: a second-order equation carries $f$, $f'$ and $f''$ in four blocks. The constraint rows of \Cref{sec:descriptor-constraints} are added to these block rows.

% CUT 2026-10-06: tab:descriptor-rows (the block rows again, a third time after
%   eq:descriptor-recurrence-rows/eq:descriptor-ode-row and the Definition's matrix);
%   its padding row now sits in the paragraph above.
% CUT 2026-10-06: "The chain at its smallest order" (f'+2f=0 as a chain, eq:descriptor-first-order).
%   The N = 4 instance of the chain is in app:cheb; the mass form below gives the k = 1 case.

\paragraph{The mass form ($k=1$).} A first-order equation or system $z'=\mathbb{A}z$ with $C$ components has in the standard form the rows $I_C\otimes\G-\mathbb{A}\otimes I_N$, one block per component. The equation already determines each derivative, so a derivative block would carry nothing new; the mass form instead multiplies the rows on the left by $I_C\otimes\Bt$ and uses $\Bt\G=\Dt$,
\begin{equation}
    (I_C\otimes\Bt)\bigl(I_C\otimes\G-\mathbb{A}\otimes I_N\bigr)
    = I_C\otimes\Dt-\mathbb{A}\otimes\Bt .
    \label{eq:mass-form}
\end{equation}
Both factors are banded and neither contains $\G$. The name is taken from the ODE literature, $\Bt$ playing the mass matrix of $\Bt\dot z=\dots$. For $f'+2f=0$ the mass form is $\Dt+2\Bt$ on $N$ unknowns, where a chain would need $2N$ unknowns and elimination through $\G$ gives the dense $\G+2I$; for a scalar equation with a constant coefficient it is the ultraspherical method's own operator. The mass form has the standard form's state, initial state and readout, and differs from it only by a banded invertible factor on the equation rows. It adds no qubits and has no derivative block to rescale (\Cref{sec:descriptor-rescaling}). Since the Carleman lifts of \Cref{sec:nonlinear} and the coupled fields of \Cref{sec:descriptor-constraints} are first-order systems, the mass form is the layout the nonlinear problems use.

Both layouts rest on one fact.

\begin{theorem}[Kernel equivalence]
\label{thm:kernel-equivalence}
Let $R_{\mathrm{std}}$ stack $\Astd$ with the constraint rows of \Cref{sec:pihm-residual}. Let $R_{\mathrm{desc}}$ stack the descriptor's recurrence, equation and padding rows with the same constraints, each reading the block of the derivative it constrains (\Cref{sec:descriptor-constraints}), and let $R_{\mathrm{mass}}$ be the mass form's stacked residual. Then
\begin{equation}
    \ker R_{\mathrm{desc}} = \bigl\{\, (p,\ \G p,\ \dots,\ \G^{k} p,\ 0) \;:\; p \in \ker R_{\mathrm{std}} \,\bigr\},
    \qquad
    \ker R_{\mathrm{mass}} = \ker R_{\mathrm{std}} .
    \label{eq:kernel-equivalence}
\end{equation}
\end{theorem}
\begin{proof}
Write $Ep=(p;\G p;\dots;\G^kp;0)$. Chain: the padding rows force the padding blocks of a kernel vector $w$ to zero, and since $\Bt$ is invertible each recurrence row gives $w_{j+1}=\Bt^{-1}\Dt\,w_j=\G w_j$, so $w=Ew_0$. The recurrence rows vanish on $Ep$ for every $p$, because $(\Bt\G-\Dt)\G^jp=0$, and every other row of $R_{\mathrm{desc}}$ acts on $Ep$ as the standard form's row of the same name acts on $p$: the equation row substitutes $\G^jp$ for $\fhat{j}$, and a constraint row reads the block that holds $\G^jp$. Hence $R_{\mathrm{desc}}Ep$ vanishes exactly when $R_{\mathrm{std}}p$ does, and $E$ is injective because its first block is the identity. Mass form: $R_{\mathrm{mass}}$ is $R_{\mathrm{std}}$ with its equation rows multiplied on the left by the invertible $I_C\otimes\Bt$ (\Cref{eq:mass-form}) and its constraint rows unchanged.
\end{proof}

The theorem concerns kernels, and it is exact. An exact kernel exists only where the solution lies in the span of the basis; otherwise both residuals have full rank, and their ground states are the smallest right singular vectors of different matrices. The descriptor measures a state's residual against a norm that carries its derivatives (\Cref{sec:descriptor-gap}), and the mass form weights the equation rows by $\Bt$. The two ground states converge to the same function as $n$ grows, at rates we measure but do not prove: on R2a at $n=7$ the decoded field's error against the analytic solution is $\res[2,sci]{build/R2a/descriptor/n7}{field_err_linf}$ in the descriptor and $\res[2,sci]{build/R2a/standard/n7}{field_err_linf}$ in the standard form (\Cref{sec:results-gap}). We therefore claim equal kernels and observe a common limit, not equal ground states.

\subsection{Constraints, Partial Differential Equations and Coupled Fields}
\label{sec:descriptor-constraints}

A data constraint sits on the block that holds the quantity it constrains. A Dirichlet condition $f(x_z)=0$ reads block~$0$ with the point-evaluation row $\Brow(x_z)$, as in the standard form. A Neumann condition $f'(x_m)=0$ reads block~$1$ with the same row, $\Brow(x_m)\fhat{1}=0$, where the standard form must write $\Brow(x_m)\G\ket{\psi}$. The regular datum of \Cref{sec:pihm-residual} reads block~$0$, so the field is decoded and scaled exactly as in the standard form.

On several axes the field is block~$0$, and each axis $a$ owns a chain $\partial_af,\dots,\partial_a^{k_a}f$ up to the highest order that the equation or a constraint takes along it. The recurrences use that axis's factors $\Bt_a,\Dt_a$, tensored with the identity on the other axes. The heat equation ($k_t=1$, $k_x=2$) carries four blocks, and Laplace's and the wave equation ($k=2$ on each axis) five, padded to eight. A mixed derivative would need a block of its own; the construction does not build one, which limits it to equations without mixed derivatives. A system of coupled first-order fields takes the mass form, one block per component. Initial and boundary \emph{functions}, as opposed to point values, enter both forms as datum slices, rows that pin a slice of the state to a known function. Datum slices are available to the standard form too, and neither form doubles the state of a linear \gls{pde}, so the descriptor's advantage on these problems is not one of state size.

\subsection{The Gap of a Carried Derivative}
\label{sec:descriptor-gap}

Carrying the derivatives costs the descriptor something the standard form does not pay. With $E$ the map of \Cref{thm:kernel-equivalence}, the descriptor's Rayleigh quotient on the states that satisfy the recurrences is
\begin{equation}
    \frac{\|R_{\mathrm{desc}}Ep\|^2}{\|Ep\|^2}
    = \frac{\|R_{\mathrm{std}}\,p\|^2}{\sum_{m=0}^{k}\|\G^{m}p\|^2},
    \label{eq:descriptor-quotient}
\end{equation}
because the recurrence rows vanish on $Ep$ and the numerator is the standard form's. By the min--max principle the descriptor's singular values are bounded above by the generalised singular values of this pair. The denominator carries every derivative, so the gap is flat in $N$ only where the equation controls each carried derivative: where no state $p$ orthogonal to the solution has a small residual against a large derivative. This is an upper bound, not a proof of flatness, and we measure the gap instead.
\todo{Try to prove that the gap is flat for a regular one-axis equation (the lower bound on the quotient of \Cref{eq:descriptor-quotient}, uniformly in $N$). If proved, put the proof in Appendix G (\Cref{app:descriptor}) and reword this paragraph, \Cref{eq:descriptor-degree} and the summary from ``measured'' to ``proved''.} The standard form measures the same states by $\|p\|$, against which their residuals are large, so its gap stays flat; its difficulty is the other side of the quotient, the norm of \Cref{eq:standard-wall}.

For a one-axis equation whose leading coefficient has no zero on $[-1,1]$, the continuum a priori estimate $\|f^{(k)}\|\le(\|Lf\|+\|\text{lower-order terms}\|)/\min|a_k|$ suggests that the quotient stays bounded away from zero on the complement of the kernel, and the measured gap is indeed flat: on R2a the descriptor's gap is $\res[3]{build/R2a/descriptor/n7}{gap}$ at $n=7$ and $\res[3]{build/R2a/descriptor/n8}{gap}$ at $n=8$. It fails in three ways.
\begin{enumerate}
    \item \emph{A leading coefficient that vanishes.} A polynomial of degree $N$ can concentrate within $\mathcal{O}(N^{-2})$ of an endpoint, where the Chebyshev points cluster, and each of its derivatives grows by a factor $\mathcal{O}(N^{2})$ there. Where $a_k$ vanishes at that endpoint, the leading term $a_kp^{(k)}$ of such a state is small, the lower-order terms are small as well, and the quotient is small against the top block. Without rescaling the gap falls as $N^{-4}$ at a simple zero at an endpoint and as $N^{-2}$ at an interior one; on the Legendre equation R3a, whose leading coefficient is $1-x^2$, it is $\res[2,sci]{build/R3a/descriptor/n8/unscaled}{gap}$ at $n=8$.
    \item \emph{Two axes.} The equation constrains only a sum of top derivatives, $f_{xx}+f_{yy}$ or $f_t-\kappa f_{xx}$. A high-frequency state with $\G_x^2p\approx-\G_y^2p$, or with its weight in the block the equation weights least, has a small residual against large blocks, and only the datum rows penalise it. On the heat equation $\partial_tf-\tfrac{1}{25}\partial_x^2f=0$ (R5b, the paper's Fig.~6d) the unscaled gap is $\res[2,sci]{build/R5b/descriptor/n5x5/unscaled}{gap}$ at $n=5$ per axis.
    \item \emph{The field block crowded out.} The ground state's weight on the field block is about $\|b\|/(\sum_m\|\G^mb\|^2)^{1/2}$, with $b$ the solution's coefficients, and the growth of the derivatives makes it small. An initial state prepared on the field block then overlaps the ground state by at most that weight, and amplification costs its reciprocal. On the stiff equation $f''+5f'+400f=0$ (R2c, the paper's Fig.~3d) the unscaled gap is $\res[2,sci]{build/R2c/descriptor/n7/unscaled}{gap}$ and the field block's weight $\res[2,sci]{build/R2c/descriptor/n7/unscaled}{field_weight}$ at $n=7$.
\end{enumerate}
The collapse is not an artefact of the lift, the padding, the arithmetic precision, the constraints, the basis or the classical reference: each was tested and ruled out, and a descriptor built in the ultraspherical basis collapses at the same rates where the leading coefficient vanishes. The tests and the full mechanism are given in \Cref{app:descriptor}.

\subsection{Block Rescaling}
\label{sec:descriptor-rescaling}

The remedy is to rescale the carried derivatives so that no block outweighs the equation's own scale. The block of order $m$ is carried divided by $\zeta^m$,
\begin{equation}
    \fhat{m} \;\longmapsto\; \zeta^{-m}\,\fhat{m},
    \qquad m = 0, \dots, k ,
    \label{eq:block-rescaling}
\end{equation}
which replaces $R_{\mathrm{desc}}$ by $R_{\mathrm{desc}}S$ with $S=\mathrm{diag}(1,\zeta,\dots,\zeta^k)\otimes I_N$. Because $S$ is diagonal and invertible, the kernel becomes $S^{-1}\ker R_{\mathrm{desc}}$ and \Cref{thm:kernel-equivalence} survives. The field block has weight one, so the readout is unchanged, and the padding keeps its weight and still decouples. What changes is the Hamiltonian, $SR^\dagger RS$. This is a congruence, not a similarity, so the spectrum moves: by \Cref{eq:descriptor-quotient} the denominator becomes $\sum_m\zeta^{-2m}\|\G^mp\|^2$, and $\zeta$ sets how much each carried derivative weighs against the residual.

In the circuit the rescaling costs nothing. Scaling the columns of block $m$ by $\zeta^m$ multiplies the coefficient of every term of the block encoding of \Cref{sec:descriptor-block-encoding} by $\zeta$ raised to the order of its column block, which adds no gate and no ancilla and changes only the subnormalisation. A general diagonal equilibration, with a different weight for each amplitude, would need a diagonal block encoding of its own at $\Theta(\dim w)$ gates; that is why the weights are constant on each block, and general equilibration appears in no circuit in this thesis.

\paragraph{The a priori rule.} $\zeta$ is computed from the equation's coefficients and the resolution, never from a solution or a spectrum. For a one-axis equation whose leading coefficient $a_k$ has no zero on $[-1,1]$,
\begin{equation}
    \zeta = 2\,\bigl(\bar a_j/\bar a_k\bigr)^{1/(k-j)},
    \label{eq:zeta-rule}
\end{equation}
with $j$ the lowest order present and $\bar a$ the mean of $|a|$ over the interval. For constant coefficients and $j=0$, $(\bar a_0/\bar a_k)^{1/k}$ is the geometric mean of the magnitudes of the characteristic roots: the scale at which a solution $e^{\lambda x}$ has $\|f^{(m)}\|\approx|\lambda|^m\|f\|$, so that dividing block $m$ by its $m$-th power balances the blocks at the equation's own scale. Where $a_k$ has a zero on the interval, $\zeta=N/2$, which makes the quotient of the concentrated states of \Cref{sec:descriptor-gap} of order one. On several axes $\zeta_a=N_a$ per axis, balanced across axes by the top-order coefficients, so that the wave equation $f_{tt}=c^2f_{xx}$ takes $(cN,N)$. The two constants, the $2$ in \Cref{eq:zeta-rule} and the $\tfrac12$ in $N/2$, were fixed once by sweeping $\zeta$ over the regular and the singular benchmark problems, and are never refitted to a problem; the rule is then checked against the sweep's best $\zeta$ on every problem. The filter's cost stays within twice its least value over a window of $\zeta$ \prov{two to eight times wide}, so $\zeta$ need only be right to about a factor of two.\todo{Appendix G: the rule's check against the grid's best $\zeta$ (within about $1.2\times$ on 27 of 30 panel-sizes, per \texttt{14\_descriptor\_gap.md} \S5.2), the window table, and the sweep that set the constants. The window's \texttt{\textbackslash prov} needs a \texttt{\textbackslash res} key or the appendix table.}

\paragraph{What it restores.} On every one-axis problem the descriptor runs, the rescaled gap no longer falls with $N$, and the field block regains its weight. The gap meant here is that of the physical subspace, the one the filter meets, since the padding alone would cap the gap of the whole spectrum at one. On R2c the gap rises from $\res[2,sci]{build/R2c/descriptor/n7/unscaled}{gap}$ to $\res[3]{build/R2c/descriptor/n7}{gap}$ and the field block's weight from $\res[2,sci]{build/R2c/descriptor/n7/unscaled}{field_weight}$ to $\res[3]{build/R2c/descriptor/n7}{field_weight}$ at $n=7$. On R3a the rescaled gap is $\res[3]{build/R3a/descriptor/n5}{gap}$ at $n=5$ and $\res[3]{build/R3a/descriptor/n8}{gap}$ at $n=8$, against $\res[2,sci]{build/R3a/descriptor/n8/unscaled}{gap}$ unscaled: it rises slowly with $n$ rather than falling.

\paragraph{What it does not remove.} Three limits remain, each a finding.
\begin{enumerate}
    \item \emph{The norm grows with $\zeta$.} Block $m$ carries $\zeta^m$, so $\|R\|\gtrsim\zeta^k\bar a_k$ and $\|H\|\gtrsim\zeta^{2k}$. With $\zeta$ constant this costs nothing. With $\zeta\propto N$, which a simple zero requires, $\|H\|=\Theta(N^{2k})$, and at a flat gap the degree at second order is about $\tilde\Theta(N^{2})$, against the standard form's $\tilde\Theta(N^4)$ (\Cref{sec:descriptor-spectrum}).
    \item \emph{Two axes.} The rescaled gap still falls with $N$, though far more slowly than unscaled: on R5b it is $\res[3]{build/R5b/descriptor/n3x3}{gap}$ at $n=3$ per axis and $\res[3]{build/R5b/descriptor/n5x5}{gap}$ at $n=5$.
    \item \emph{A double zero is beyond any block rescaling.} If $a_k$ vanishes to second order at an endpoint, the concentrated states have quotient $\mathcal{O}(N^{-8}\zeta^{2k})$, which is the unscaled $N^{-8}$ at $\zeta=1$, so balancing them needs $\zeta\propto N^{4/k}$. At $k=2$ that is $\zeta\propto N^2$, at which $\|H\|=\Theta(N^8)$, the standard form's norm. The paper's own Legendre panels of Figs.~4c and 4d (R3c and R3d) multiply the associated Legendre operator by $1-x^2$ and so have exactly this double zero. The descriptor instead runs the equivalent problem for $v=f/\sqrt{1-x^2}$, whose leading coefficient has simple zeros, and we make no claim for the double zero.
\end{enumerate}
Block rescaling is the descriptor's counterpart of the diagonal preconditioner of the ultraspherical method \cite{olver2013}, which bounds the condition number of its banded system. Here the diagonal scales the state's columns and is constant on each block, which is why it costs nothing in the circuit.

\subsection{The Block Encoding: One Flat Linear Combination}
\label{sec:descriptor-block-encoding}

Every block of $\Adesc$ is a short sum of \emph{bands}, each a diagonal followed by a truncated shift. The diagonal of $\Dt$ is $2c$ on column $c$, affine in the index except at column~1, where the normalisation of $T_0$ intervenes; the diagonals of $\Bt$ and of a constant multiplication are constant except at their ends; and a variable coefficient's multiplication and the Galerkin truncation contribute a few further exceptions. An affine diagonal is a sum of $n+1$ Pauli-$Z$ terms, and each exception is one reflection $Q_m$:
\begin{equation}
    \operatorname{diag}(c)_{c=0}^{N-1} = \sum_{j=0}^{n-1} 2^{j}\,\frac{I - Z_j}{2},
    \qquad
    Q_m = I - 2\ket{m}\bra{m} .
    \label{eq:diag-index}
\end{equation}
$\Adesc$ is therefore one flat \gls{lcu} of $\mathcal{O}(n)$ terms,
\begin{equation}
    \Adesc = \sum_{t=1}^{L} c_t\,U_t,
    \qquad
    U_t = (\text{field word}) \otimes \mathrm{SHIFT}^{s_t}\,P_t,
    \qquad
    \alpha_A = \sum_{t=1}^{L} |c_t|,
    \label{eq:descriptor-lcu}
\end{equation}
where the field word acts on the block register, $\mathrm{SHIFT}^{s_t}$ is the truncated shift by $s_t$, and $P_t$ is a $Z$-string or a reflection. An exactly affine diagonal whose end values share a sign is encoded tightly, at $\alpha=\max_c|a+bc|$, so the dominant term $\Dt$ is encoded at its own norm $2(N-1)$ and its exceptions cost $\mathcal{O}(1)$. A truncated shift of the $N$-dimensional register is not unitary; it is applied on the register widened by one overflow qubit, so that rows falling outside the interval land in the half that the final projection discards. That qubit is the shift's only ancilla.

\Cref{alg:ir-to-lcu} compiles the terms into the circuit of \Cref{fig:descriptor-lcu}.

\placeholderfigure{\textbf{The whole descriptor residual is one PREPARE, one SELECT and one PREPARE$^\dagger$.} PREPARE loads $\sqrt{|c_t|/\alpha_A}$ over the joint index of the factored selector registers. SELECT applies each term's field word on the block register ($n_f$ qubits), then on the system register ($n$ qubits, widened by one overflow qubit) its diagonal, a $Z$-string or a reflection, and its truncated shift. Post-selecting the selectors, the overflow qubit and the flag on $\ket{0}$ leaves $\Adesc/\alpha_A$, with the block rescaling carried entirely in the coefficients $c_t$.}{fig:descriptor-lcu}{Quantikz circuit: selector registers (PREPARE loads $\sqrt{|c_t|/\alpha_A}$ over their joint index), field register ($n_f$ qubits), system register ($n$ qubits) widened by one overflow qubit; SELECT applies each term's field word, diagonal ($Z$-string or reflection) and truncated shift (a ladder of multi-controlled-$X$ gates); PREPARE$^\dagger$; post-selection on the selectors, the overflow and the flag.}

\begin{algorithm}[htbp]
    \caption{Compiling the rows of $\Adesc$ into PREPARE and SELECT.}
    \label{alg:ir-to-lcu}
    \KwIn{the blocks of $\Adesc$ as bands, the block weights $\zeta$}
    \KwOut{PREPARE and SELECT circuits; the subnormalisation $\alpha_A$}
    Write every block as a sum of bands: a diagonal and a truncated shift\;
    Fit each diagonal as affine in the column index plus a few exceptions, \Cref{eq:diag-index}\;
    Merge terms with equal axis factors across blocks; expand the resulting block-register pattern in Pauli words\;
    Multiply the coefficient of each term by $\zeta^{m}$ for its column block $m$\;
    $\alpha_A \leftarrow \sum_{t} |c_t|$; PREPARE loads $\operatorname{sgn}(c_t)\sqrt{|c_t|/\alpha_A}$ on the joint selector index of term $t$\;
    SELECT: controlled on each selector register, apply the field word, then on each axis the diagonal and the shift\;
    \Return{PREPARE, SELECT, $\alpha_A$}\;
\end{algorithm}

% CUT 2026-10-06: "Inputs and outputs" (restated the algorithm's KwIn/KwOut).
% MOVED 2026-10-06 to 6_nonlinear.tex (the Carleman coupling's matchings, where the lift is built):
%   a large irregular coupling is split into matchings, each with at most one entry in every row and
%   column, and a greedy colouring uses at most 2*Delta - 1 of them, Delta the largest row degree.

\paragraph{The terms (lines 1--4).} Lines 1 and 2 are the decomposition above. Line 3 gathers terms whose axis factors agree across the rows and columns they join, so that the block register carries their pattern once, as a small matrix expanded in Pauli words; the recurrences' common $\Dt$ is paid for once and not $k$ times. Line 4 is the whole of the block rescaling.

\paragraph{PREPARE and SELECT (lines 5--6).} The selectors are factored: one register for the field operation and, on each axis, one for the diagonal and one for the shift, each only as wide as the number of distinct values its terms use. SELECT then applies each distinct operation once, controlled on its own register, rather than once per term, which is what keeps the ancilla at $\mathcal{O}(\log n)$ (\Cref{sec:descriptor-cost}). The rank-one rows, the sources and the data constraints, are encoded through the feature map, with a rotation tree loading a source's coefficients, so no dense state preparation appears anywhere. The constructions and the term-by-term accounting are given in \Cref{app:term-ir}.

\subsection{Cost of the Descriptor Encoding}
\label{sec:descriptor-cost}

\begin{proposition}[Cost of the descriptor encoding]
\label{prop:descriptor-cost}
For a descriptor residual of a fixed problem, the flat \gls{lcu} of \Cref{eq:descriptor-lcu} is an exact block encoding of the system rows at $\alpha_A=\sum_t|c_t|$, with $\mathcal{O}(\log n)$ ancilla and $\mathcal{O}(n^2)$ gates. For the second-order family R2,
\begin{equation}
    \lceil \log_2(n+4) \rceil + 8\ \text{ancilla},
    \qquad
    \mathcal{O}(n^2)\ \text{gates} .
    \label{eq:descriptor-cost}
\end{equation}
On a regular one-axis equation at fixed $\zeta$, the subnormalisation of the system rows is $\alpha_A/\|\Adesc\|_2=1+\mathcal{O}(1/N)$.
\end{proposition}

The ancilla are the selector registers, the overflow qubit and the flags that project the field terms. Only the diagonal selector's width depends on $n$: for R2 it indexes the identity, the $n$ gates $Z_j$ and three reflections, so it has $\lceil\log_2(n+4)\rceil$ qubits, while the field and shift selectors have widths fixed by the problem. The gate count follows because each distinct operation is applied once and the shift's multi-controlled ladder costs $\mathcal{O}(n^2)$. The proof is given in \Cref{app:proof-descriptor-cost}, and the ancilla count is verified for R2's system rows at $n=2$ to $12$.\todo{Appendix G (\texttt{app:proof-descriptor-cost}) is still a scaffold: write the ancilla and gate count and the $1+\mathcal{O}(1/N)$ ratio (the dominant term $\zeta^j\Dt$ is encoded at its norm), then check this proposition against it. The ratio is tested, not proved, until then.}

The ancilla count is of the least order among exact encodings of this kind. A combination whose SELECT indexes $L$ distinct operations needs $\lceil\log_2L\rceil$ selector qubits \cite{chakraborty2025qsvt}, and an affine diagonal has exactly $n+1$ non-zero Walsh coefficients, so any \gls{lcu} that expands its diagonals in Pauli strings needs $\Omega(\log n)$ of them. We do not claim optimality against every encoding. The structured encoding of the standard form needs $n-1$ selector qubits for $\G$ alone, because $\Uodd$ is a sum of $N/2$ shifts (\Cref{sec:standard-cost}), so its ancilla grow as $\Omega(n)$ where the descriptor's grow as $\log n$: the difference between a dense operator and a banded one. The better scaling is not yet visible at the sizes that circuits reach. For the whole reflection on R2a the descriptor needs $\res{encode/R2a/descriptor/n2}{n_anc_qubits}$ ancilla at $n=2$ and $\res{encode/R2a/descriptor/n5/verify=probe}{n_anc_qubits}$ at $n=5$, against $\res{encode/R2a/standard/n2}{n_anc_qubits}$ and $\res{encode/R2a/standard/n5/verify=probe}{n_anc_qubits}$ for the standard form.

The ratio $\alpha/\|\cdot\|$ is the factor by which an encoding exceeds the least possible subnormalisation, and the filter's degree is linear in it. The proposition concerns the system rows; for the whole stacked residual, constraints included, the ratio is measured. It approaches one on regular one-axis problems, $\alpha_R/\|R\|_2=\res[3]{looseness/stateprep/R2a/descriptor/n8/T3/filter=minimax}{alpha_ratio}$ on R2a at $n=8$ (a resource estimate). It stays near two where the leading coefficient vanishes, $\res[3]{looseness/stateprep/R3a/descriptor/n7/T3/filter=minimax}{alpha_ratio}$ on R3a at $n=7$, because $\zeta\Dt$, $\zeta^2\Bt$ and $\zeta^2a_2$ are then comparable and comparable terms add in $\ell^1$ even when each is tight. On two axes it is $\res[3]{looseness/stateprep/R5b/descriptor/n4x4/T3/filter=minimax}{alpha_ratio}$ on R5b at $n=4$ per axis. The limit of one therefore holds for regular one-axis problems only. For comparison, the structured standard-form encoding is at $\res[3]{looseness/stateprep/R2a/standard/n8/T3/filter=minimax}{alpha_ratio}$ on R2a and $\res[2,sci]{looseness/stateprep/R3a/standard/n7/T3/filter=minimax}{alpha_ratio}$ on R3a, the latter the lifted multiplier's looseness of \Cref{sec:standard-hamiltonian}.

The gates per query go the other way on the constant-coefficient problems. On R2a, at the sizes reached, one query of the descriptor costs more than one of the structured standard form, having more terms and a block register: the transpiled two-qubit count of the reflection is $\res{encode/R2a/descriptor/n2}{cx_count}$ against $\res{encode/R2a/standard/n2}{cx_count}$ at $n=2$ and $\res{encode/R2a/descriptor/n5/verify=probe}{cx_count}$ against $\res{encode/R2a/standard/n5/verify=probe}{cx_count}$ at $n=5$. The ratio falls with $n$, but both counts are $\mathcal{O}(n^2)$, so the data do not show that it crosses one. The descriptor's advantage is therefore the number of queries, the filter's degree, and not the gates per query, and its chain state is larger by the block register. The next section is where the advantage lies.
% REVISE (conflict with CLAUDE.md Standing Facts, "at reachable n the descriptor's gates
%   per query exceed the structured standard form's"): pihm/docs/13_descriptor.md §3 shows
%   the descriptor's reflection CHEAPER than the structured standard form's on the lifted
%   panels from n = 4 (R3a: 38864 vs 43108 cx at n = 4, 48542 vs 52038 at n = 5; R4a from
%   n = 3), and dearer on R1, R2, R4b-d and R5. The sentence above is therefore scoped to
%   R2a. Decide whether the Standing Fact should be scoped to constant coefficients and,
%   if so, quote the R3a counts here by key (encode/R3a/{standard,descriptor}/n5/verify=probe).

\subsection{Consequences for the Spectrum and the Filter}
\label{sec:descriptor-spectrum}

Every block of the descriptor is $\mathcal{O}(N)$ at fixed $\zeta$, the recurrences' $\zeta^j\Dt$ being the largest, so $\|R_{\mathrm{desc}}S\|=\Theta(N)$ and $\|\Hdesc\|=\Theta(N^2)$ whatever the order $k$: the derivative is banded, and its norm is that of $\Dt$ rather than of $\G^k$. With the encoding at the norm, $\alpha_H/\|H\|$ falls from $\res[3]{looseness/stateprep/R2a/standard/n8/T3/filter=minimax}{hamiltonian_ratio}$ for the standard form on R2a at $n=8$ to $\res[3]{looseness/stateprep/R2a/descriptor/n8/T3/filter=minimax}{hamiltonian_ratio}$ for the descriptor. Since the degree is $d=\Theta(\sqrt{\alpha_H/\Delta}\log1/\varepsilon)$ (\Cref{eq:degree}), the rescaled gap, which is flat on a one-axis \gls{ode} whose leading coefficient does not vanish, gives for that class of problems
\begin{equation}
    \|\Hdesc\| = \Theta(N^2),
    \qquad
    \Delta = \Theta(1)\ \text{(measured)}
    \quad\Longrightarrow\quad
    d = \tilde\Theta(N),
    \label{eq:descriptor-degree}
\end{equation}
against $\tilde\Theta(N^{2k})$ for the standard form with the best exact encoding (\Cref{eq:standard-wall}). Where the leading coefficient vanishes, or on two axes, the rescaling's weights raise the norm and the degree is about $\tilde\Theta(N^2)$ or faster; for a double zero it is the standard form's or worse (\Cref{tab:descriptor-regimes}). Neither form is polylogarithmic in $N$: $\tilde\Theta(N)$ is polynomial in $N$ and so exponential in $n$, and the claim is a lower exponent, not an exponential speed-up. The degree is that of the minimax edge filter, which prepares every comparison in this thesis; for polynomial filters in general the floor is no stronger than Lin and Tong's optimality supports \cite{lintong2020groundstate}.

\begin{table}[htbp]
    \centering
    \caption{\textbf{The descriptor lowers the degree's exponent where the equation controls every carried derivative.} Norm, gap and filter degree of each form by class of problem, with the block weight $\zeta$ set by the a priori rule of \Cref{sec:descriptor-rescaling}. Norms follow from the block weights and are proved; gaps are measured.}
    \label{tab:descriptor-regimes}
    \small
    \begin{tabular}{@{}l l l l l@{}}
        \toprule
        Case & $\zeta$ & $\|H\|$ & Gap & Degree \\
        \midrule
        Standard, order $k$, constant coefficients & -- & $\Theta(N^{4k})$ & flat & $\tilde\Theta(N^{2k})$ \\
        Descriptor, one axis, $a_k \ne 0$ on $[-1,1]$ & fixed & $\Theta(N^2)$ & flat & $\tilde\Theta(N)$ \\
        Descriptor, one axis, simple zero, $k=2$ & $N/2$ & $\Theta(N^4)$ & flat & $\tilde\Theta(N^2)$ \\
        Descriptor, two axes & $N_a$ per axis & $\Theta(N^4)$ & falls & above $\tilde\Theta(N^2)$ \\
        Descriptor, double zero & $\propto N^2$ & $\Theta(N^8)$ & -- & standard's or worse \\
        \bottomrule
    \end{tabular}
\end{table}

% MOVED 2026-10-06 to 8_results.tex, sec:results-scaling (the measured degree exponents between n = 7 and
%   n = 8: R2a descriptor stateprep/R2a/descriptor/growth/T3/filter=minimax vs standard .../standard/growth/...;
%   R3a descriptor .../R3a/descriptor/growth/...; R5b descriptor and standard .../R5b/*/growth/...). They are
%   two-point slopes of a degree with logarithmic factors, so say "consistent with", not "bears out".
On R2a at $n=8$ the estimated degree is $\res[2,sci]{stateprep/R2a/descriptor/n8/T3/filter=minimax}{degree}$ for the descriptor, against $\res[2,sci]{stateprep/R2a/standard/n8/T3/filter=minimax}{degree}$ for the standard form as built and $\res[2,sci]{stateprep/R2a/standard/n8/T3/filter=minimax}{ideal_degree}$ at its ideal-encoding bound, all three resource estimates; the measured growth with $N$ is reported in \Cref{sec:results-scaling}. Every comparison is quoted both as built and at $\alpha=\|R\|_2$, which contains neither form's encoding, and where the order reverses, as it does on R3 at the ideal bound and on the heat equation at small $n$, the reversal is reported as a finding (\Cref{sec:results-scaling,sec:results-gap}).

Two further facts bound what the comparison says. The first is the relative gap: on R2a at $n=7$ it is $\Delta/\|\Hstd\|=\res[2,sci]{build/R2a/standard/n7}{relative_gap}$ for the standard form and $\res[2,sci]{build/R2a/descriptor/n7}{relative_gap}$ for the descriptor. An eigensolver applied to $\Hstd$ resolves its ground state only to about machine precision divided by that ratio, and at these sizes it fails, whereas the singular value decomposition of $R$, whose relative scale is $\Theta(N^{-2k})$ rather than $\Theta(N^{-4k})$, succeeds; every classical reference in this thesis is therefore that decomposition, and the consequence of the small relative gap is the degree, not an unresolvable ground state. The second concerns a single query: one post-selected application of a block encoding of $R$ to a Haar-random state of dimension $D$ succeeds with mean probability $\|R\|_F^2/(D\alpha_R^2)$, which is $\Theta(1)$ for the descriptor and $\Theta(N^{-1})$ per axis for the structured standard form, whose dense $\G^k$ dominates the Frobenius norm. This compares the encodings only; it is not a cost of the preparation, which post-selects no single query.

\subsection{Summary}
\label{sec:descriptor-summary}

The residual can be reformulated so that its Hamiltonian is banded without changing its solution, and the reformulation has a price.
\begin{enumerate}
    \item Keeping the banded factors $\Bt$ and $\Dt$ of $\G$ and carrying the derivatives as unknowns, as a chain for equations of order two and above or in mass form for first-order equations and systems, makes the residual banded with its kernel exactly the standard form's (\Cref{thm:kernel-equivalence}).
    \item Its block encoding is one flat \gls{lcu} with $\lceil\log_2(n+4)\rceil+8$ ancilla for R2's system rows, of the least order among exact \gls{lcu} encodings, and its subnormalisation tends to the norm on regular one-axis problems (\Cref{prop:descriptor-cost}). Its price is a larger state, by the chain's block register, and, on the constant-coefficient problems, more gates per query at the sizes reached.
    \item The Hamiltonian's norm is $\Theta(N^2)$, and for a one-axis equation whose leading coefficient does not vanish the filter's degree is $\tilde\Theta(N)$, against the standard form's $\tilde\Theta(N^{2k})$ (\Cref{eq:descriptor-degree}). Neither is polylogarithmic.
    \item \emph{The limit.} The carried derivatives erode the gap where the leading coefficient vanishes and on two axes. An a priori block rescaling restores it on one axis at no gate cost, at the price of a larger norm where the leading coefficient vanishes; on two axes the gap still falls, and a double zero is beyond any block rescaling.
\end{enumerate}
\Cref{sec:results-scaling,sec:results-gap} test the cost and the gap on the benchmark problems. In one sentence: carrying the intermediate derivatives instead of eliminating them turns the physics-informed residual from dense to banded with its kernel unchanged, so that the Hamiltonian's norm falls from $\Theta(N^{4k})$ to $\Theta(N^2)$; the price is a block register and a gap that the derivative blocks erode, which an a priori rescaling restores on one axis at no gate cost.

% CHECKLIST (marker)
%   1. Could a referee re-derive A_desc from the body plus Appendices B and G?
%   2. Is the Hilbert-space and per-query cost volunteered alongside the gains?
%   3. Are ultraspherical and first-order least squares credited at sec:bg-spectral
%      and sec:descriptor-idea/sec:descriptor-construction?
%   4. Is the cost result a proposition AND interpreted?
%   5. Is the ancilla claim O(log n), with optimality hedged to exact LCU?
%   6. Is the gap's failure owned with its mechanism (sec:descriptor-gap)?
%   7. Is the rescaling's limits stated (two axes; a double zero)?
%   8. Does every body equation pass the equation test?
% TRAPS  "O(1) ancilla"; "+ 6 ancilla" (now + 8); "bit-identical a-block"; "the
%        gap is Theta(1)" unscoped; "rescaling costs Theta(n_f) gates" (it is a
%        priori, no gates, no ancilla); "zeta fitted"; "only the descriptor can
%        prepare nonlinear solutions".
